% 14-generics.tex — Clause 14: Generics and compile-time evaluation
\chapter{Generics and compile-time evaluation}
\label{sec:14-generics}
\index{generics}
\index{compile-time evaluation}
\index{comptime}

\section{Overview}
\label{sec:gen.overview}

\pnum
Lain implements genericity through compile-time function evaluation (CTFE).
Generic types and functions accept \textit{type parameters} (declared as
\lain{name type}) or \textit{compile-time value parameters} (declared as
\lain{comptime name non-type}).  The compiler evaluates them at compile time
and monomorphizes the result.  There is no type erasure and no runtime
polymorphism.

\section{Type parameters}
\label{sec:gen.type-params}
\index{type parameter}

\pnum
A parameter declared \lain{name type} is a type parameter.  Types are
implicitly compile-time — the \lain{comptime} keyword shall \emph{not}
precede a type parameter.  At every call site, the argument shall be a type
expression known at compile time.

\begin{laincode}
func identity(T type, x T) T {
    return x
}

var n = identity(int, 42)       // T = int, monomorphized
var s = identity(u8[:0], "hi")  // T = u8[:0], separate monomorphization
\end{laincode}

\pnum
Type parameters are evaluated entirely at compile time; they produce
no runtime code.

\begin{constraint}
  The \lain{comptime} keyword shall not appear before a parameter whose
  type is \lain{type}.  \lain{func f(comptime T type)} is ill-formed;
  write \lain{func f(T type)} instead.
\end{constraint}

\section{Compile-time value parameters}
\label{sec:gen.comptime-params}
\index{comptime parameter}

\pnum
A parameter declared \lain{comptime name non-type} is a compile-time value
parameter.  The argument shall be a constant expression of the declared type
known at compile time.

\begin{laincode}
type RingBuffer(T type, comptime N int) = type {
    items T[N]
    head  usize
    tail  usize
}
\end{laincode}

\begin{constraint}
  A compile-time evaluation context (CTFE) shall not call a \lain{proc}
  or an \lain{extern proc}; only \lain{func} calls are admissible. A
  violation is \illformed{}~(\diagcode{E101}).
\end{constraint}

\section{Type as value}
\label{sec:gen.type-value}
\index{type as value}

\pnum
A \lain{func} with return type \lain{type} is a type constructor.  Its
body shall consist only of expressions and statements that are evaluable
at compile time.  The return value is a type that the caller may use in
type annotations and variable declarations.

\begin{laincode}
type Option(T type) = type {
    Some { value T }
    None
}

var x = Option(int).Some(42)
\end{laincode}

\pnum
Type constructors may also be written as \lain{func} declarations returning
\lain{type}:

\begin{laincode}
func Option(T type) type {
    return type {
        Some { value T }
        None
    }
}
\end{laincode}

\begin{constraint}
  A function returning \lain{type} shall be declared as \lain{func} (not
  \lain{proc}).
\end{constraint}

\section{Monomorphization}
\label{sec:gen.monomorphization}
\index{monomorphization}

\pnum
For each unique combination of comptime arguments at call sites, the
implementation produces a separate instantiation of the function.  The
name of each instantiation in the translation output is
\idbehav{determined by the implementation}.

\pnum
Monomorphization is a compile-time-only operation; at runtime, each
call site calls a specific, non-generic function.

\section{Comptime expression blocks}
\label{sec:gen.comptime-block}
\index{comptime block}

\pnum
A \lain{comptime \{ expr \}} expression evaluates \textit{expr} at compile
time.  The result is a compile-time constant that may be used as a type
annotation, array size, or parameter constraint value.

\begin{constraint}
  A \lain{comptime} expression shall be pure and terminating (i.e., it may
  only use \lain{func} calls and the constructs permitted in \lain{func}
  bodies).
\end{constraint}

\section{Comptime if}
\label{sec:gen.comptime-if}
\index{comptime if}

\pnum
\lain{comptime if cond \{ ... \} else \{ ... \}} is a compile-time
conditional.  Only the taken branch is type-checked and emitted; the other
branch is entirely discarded.  \textit{cond} shall evaluate to a
compile-time boolean constant.

\begin{note}
  \lain{comptime if} enables platform-specific code selection without
  dead-code warnings.  It is the preferred way to write code that
  varies based on compile-time parameters.
\end{note}

\section{Comptime recursion depth}
\label{sec:gen.recursion-depth}
\index{comptime recursion depth}

\pnum
The maximum depth of comptime function evaluation is
\idbehav{the implementation's comptime evaluation depth limit; the
reference implementation enforces a limit of 256}.  A comptime evaluation
that exceeds this depth is \illformed{}~(\diagcode{E014}).

\section{Mode preservation in generics}
\label{sec:gen.mode-preserve}
\index{mode preservation}

\pnum
When a type parameter \lain{T} is substituted with a concrete type, the
ownership mode of the parameter is preserved.  If a parameter was declared
\lain{var T}, and \lain{T} is substituted with a type that has mode
\lain{shared} in its declaration, the substituted parameter retains the
\lain{var} mode.

\begin{note}
  This rule ensures that mode annotations in generic functions behave
  consistently regardless of the concrete type argument.
\end{note}
